\FloatBarrier
\appendix
\section{Thời gian thành phần đầy đủ}
\label{app:phases}
Các bảng sau tách build/load/solve để đọc cùng work và total trong phần chính. Đơn vị là giây. Các giá trị rất nhỏ được giữ thêm chữ số để không bị hiển thị thành 0; điều này không hàm ý phép đo có độ chính xác tuyệt đối tới chữ số cuối.
\begin{table}[pos=H]
\centering
\small
\caption{Thời gian sinh CNF SAT (giây).}
\label{tab:sat_build}
\setlength{\tabcolsep}{4pt}
\begin{tabular}{@{}lrrrrr@{}}
\toprule
$N$ & Binomial & Binary & Sequential & Commander & Product \\
\midrule
4 & 0.000101 & 0.000145 & 0.000145 & 0.000219 & 0.000120 \\
8 & 0.000337 & 0.000425 & 0.000333 & 0.000560 & 0.000584 \\
9 & 0.000436 & 0.000571 & 0.000449 & 0.000743 & 0.000670 \\
16 & 0.004802 & 0.003365 & 0.002613 & 0.004080 & 0.003813 \\
25 & 0.009329 & 0.006029 & 0.004446 & 0.007701 & 0.006155 \\
27 & 0.0113 & 0.006435 & 0.004808 & 0.009049 & 0.007402 \\
36 & 0.0268 & 0.0121 & 0.006959 & 0.0133 & 0.0105 \\
49 & 0.0617 & 0.0252 & 0.0125 & 0.0232 & 0.0188 \\
64 & 0.1360 & 0.0372 & 0.0197 & 0.0388 & 0.0279 \\
81 & 0.2839 & 0.0639 & 0.0270 & 0.0570 & 0.0376 \\
100 & 0.5188 & 0.0995 & 0.0432 & 0.0918 & 0.0523 \\
125 & 1.121 & 0.1577 & 0.0598 & 0.1373 & 1S+2TO \\
216 & 5.943 & 0.5317 & TO & 0.3939 & TO \\
343 & TO & 1.532 & TO & TO & TO \\
512 & TO & 3.562 & TO & TO & TO \\
729 & TO & TO & TO & TO & TO \\
1000 & TO & TO & TO & TO & TO \\
1331 & TO & TO & TO & TO & TO \\
\bottomrule
\end{tabular}
\par\smallskip\begin{minipage}{\linewidth}\footnotesize Giá trị số: trung vị của ba lượt SAT có nghiệm hợp lệ. TO: ba TIMEOUT; 2TO+1U: hai TIMEOUT và một UNKNOWN; 1S+2TO: một SAT và hai TIMEOUT. Không thay lượt chưa hoàn tất bằng 0 hoặc 100 giây.\end{minipage}
\end{table}

\begin{table}[pos=H]
\centering
\small
\caption{Thời gian nạp CNF SAT (giây).}
\label{tab:sat_load}
\setlength{\tabcolsep}{4pt}
\begin{tabular}{@{}lrrrrr@{}}
\toprule
$N$ & Binomial & Binary & Sequential & Commander & Product \\
\midrule
4 & 0.000136 & 0.000187 & 0.000175 & 0.000203 & 0.000136 \\
8 & 0.000739 & 0.001215 & 0.000655 & 0.000703 & 0.000750 \\
9 & 0.000943 & 0.001127 & 0.000841 & 0.000904 & 0.001134 \\
16 & 0.008330 & 0.003618 & 0.002731 & 0.003581 & 0.002762 \\
25 & 0.0191 & 0.0110 & 0.007011 & 0.008278 & 0.008171 \\
27 & 0.0243 & 0.0128 & 0.007763 & 0.0106 & 0.009367 \\
36 & 0.0590 & 0.0246 & 0.0140 & 0.0179 & 0.0155 \\
49 & 0.1459 & 0.0474 & 0.0274 & 0.0339 & 0.0337 \\
64 & 0.3415 & 0.0863 & 0.0448 & 0.0556 & 0.0431 \\
81 & 0.6932 & 0.1495 & 0.0762 & 0.0893 & 0.0701 \\
100 & 1.315 & 0.2409 & 0.1190 & 0.1438 & 0.1045 \\
125 & 2.732 & 0.3835 & 0.1767 & 0.2278 & 1S+2TO \\
216 & 15.416 & 1.298 & TO & 0.7232 & TO \\
343 & TO & 4.263 & TO & TO & TO \\
512 & TO & 9.401 & TO & TO & TO \\
729 & TO & TO & TO & TO & TO \\
1000 & TO & TO & TO & TO & TO \\
1331 & TO & TO & TO & TO & TO \\
\bottomrule
\end{tabular}
\par\smallskip\begin{minipage}{\linewidth}\footnotesize Giá trị số: trung vị của ba lượt SAT có nghiệm hợp lệ. TO: ba TIMEOUT; 2TO+1U: hai TIMEOUT và một UNKNOWN; 1S+2TO: một SAT và hai TIMEOUT. Không thay lượt chưa hoàn tất bằng 0 hoặc 100 giây.\end{minipage}
\end{table}

\begin{table}[pos=H]
\centering
\small
\caption{Thời gian quanh lời gọi giải SAT (giây).}
\label{tab:sat_solve}
\setlength{\tabcolsep}{4pt}
\begin{tabular}{@{}lrrrrr@{}}
\toprule
$N$ & Binomial & Binary & Sequential & Commander & Product \\
\midrule
4 & 0.000061 & 0.000085 & 0.000064 & 0.000079 & 0.000063 \\
8 & 0.000265 & 0.000334 & 0.000212 & 0.000641 & 0.000670 \\
9 & 0.000154 & 0.000459 & 0.000261 & 0.000557 & 0.000863 \\
16 & 0.0178 & 0.0111 & 0.0110 & 0.0702 & 0.007254 \\
25 & 0.005146 & 0.0139 & 0.0421 & 0.0209 & 0.0211 \\
27 & 0.0163 & 0.007481 & 0.0906 & 0.0446 & 0.0832 \\
36 & 0.008011 & 0.0567 & 0.1181 & 0.0870 & 0.0740 \\
49 & 0.006487 & 0.1920 & 0.5517 & 0.0614 & 0.7420 \\
64 & 0.0229 & 0.1026 & 3.839 & 0.2145 & 18.884 \\
81 & 0.1186 & 0.1314 & 2.839 & 0.2276 & 3.541 \\
100 & 0.0554 & 0.1014 & 4.554 & 1.158 & 17.169 \\
125 & 0.0874 & 0.1873 & 16.062 & 2.650 & 1S+2TO \\
216 & 1.237 & 0.4549 & TO & 2.818 & TO \\
343 & TO & 2.300 & TO & TO & TO \\
512 & TO & 66.373 & TO & TO & TO \\
729 & TO & TO & TO & TO & TO \\
1000 & TO & TO & TO & TO & TO \\
1331 & TO & TO & TO & TO & TO \\
\bottomrule
\end{tabular}
\par\smallskip\begin{minipage}{\linewidth}\footnotesize Giá trị số: trung vị của ba lượt SAT có nghiệm hợp lệ. TO: ba TIMEOUT; 2TO+1U: hai TIMEOUT và một UNKNOWN; 1S+2TO: một SAT và hai TIMEOUT. Không thay lượt chưa hoàn tất bằng 0 hoặc 100 giây.\end{minipage}
\end{table}

\begin{table}[pos=H]
\centering
\small
\caption{Thời gian xây mô hình CP/ILP (giây).}
\label{tab:native_build}
\setlength{\tabcolsep}{4pt}
\begin{tabular}{@{}lrrrr@{}}
\toprule
$N$ & CP-SAT & CPLEX CP & CPLEX MIP & Gurobi \\
\midrule
4 & 0.001613 & 0.001550 & 0.003312 & 0.002651 \\
8 & 0.001404 & 0.001431 & 0.004678 & 0.003296 \\
9 & 0.002194 & 0.001688 & 0.004702 & 0.003258 \\
16 & 0.001679 & 0.001360 & 0.008655 & 0.007264 \\
25 & 0.001928 & 0.001458 & 0.0151 & 0.0136 \\
27 & 0.002086 & 0.001694 & 0.0155 & 0.0140 \\
36 & 0.002265 & 0.001576 & 0.0240 & 0.0241 \\
49 & 0.002977 & 0.001628 & 0.0376 & 0.0424 \\
64 & 0.003134 & 0.001833 & 0.0629 & 0.0662 \\
81 & 0.003520 & 0.002371 & 0.0812 & 0.1047 \\
100 & 0.0114 & 0.002150 & 0.1203 & 0.1558 \\
125 & 0.004471 & 0.002480 & 0.2021 & 0.2309 \\
216 & 0.006869 & 0.003330 & 0.6064 & 0.7344 \\
343 & 2TO+1U & 0.004576 & 1.529 & 1.877 \\
512 & TO & 0.006161 & 3.817 & 4.462 \\
729 & 2TO+1U & 0.008833 & TO & 10.068 \\
1000 & TO & 0.0124 & 16.685 & 20.251 \\
1331 & TO & 0.0168 & TO & TO \\
\bottomrule
\end{tabular}
\par\smallskip\begin{minipage}{\linewidth}\footnotesize Giá trị số: trung vị của ba lượt SAT có nghiệm hợp lệ. TO: ba TIMEOUT; 2TO+1U: hai TIMEOUT và một UNKNOWN; 1S+2TO: một SAT và hai TIMEOUT. Không thay lượt chưa hoàn tất bằng 0 hoặc 100 giây.\end{minipage}
\end{table}

\begin{table}[pos=H]
\centering
\small
\caption{Thời gian quanh lời gọi giải CP/ILP (giây).}
\label{tab:native_solve}
\setlength{\tabcolsep}{4pt}
\begin{tabular}{@{}lrrrr@{}}
\toprule
$N$ & CP-SAT & CPLEX CP & CPLEX MIP & Gurobi \\
\midrule
4 & 0.0109 & 0.0678 & 0.008737 & 0.007854 \\
8 & 0.0203 & 0.0629 & 0.007844 & 0.006175 \\
9 & 0.0222 & 0.0604 & 0.008861 & 0.005898 \\
16 & 0.0442 & 0.0556 & 0.0130 & 0.0320 \\
25 & 0.0943 & 0.0641 & 0.0190 & 0.006594 \\
27 & 0.1091 & 0.0594 & 0.0197 & 0.006246 \\
36 & 0.2132 & 0.0605 & 0.0298 & 0.007075 \\
49 & 0.4231 & 0.0602 & 0.0457 & 0.1496 \\
64 & 0.7286 & 0.0696 & 0.0770 & 0.009394 \\
81 & 1.392 & 0.0714 & 0.1084 & 0.2311 \\
100 & 3.342 & 0.0706 & 0.1740 & 0.2749 \\
125 & 4.953 & 0.0707 & 0.2679 & 0.9562 \\
216 & 29.904 & 0.0830 & 1.150 & 1.936 \\
343 & 2TO+1U & 0.0994 & 3.436 & 0.0926 \\
512 & TO & 0.1766 & 7.089 & 0.2577 \\
729 & 2TO+1U & 0.2066 & TO & 0.6932 \\
1000 & TO & 0.4581 & 55.126 & 1.616 \\
1331 & TO & 0.5043 & TO & TO \\
\bottomrule
\end{tabular}
\par\smallskip\begin{minipage}{\linewidth}\footnotesize Giá trị số: trung vị của ba lượt SAT có nghiệm hợp lệ. TO: ba TIMEOUT; 2TO+1U: hai TIMEOUT và một UNKNOWN; 1S+2TO: một SAT và hai TIMEOUT. Không thay lượt chưa hoàn tất bằng 0 hoặc 100 giây.\end{minipage}
\end{table}

\FloatBarrier
\section{Độ phân tán và cách tính}
\label{app:dispersion}
\begin{table}[pos=H]
\centering
\footnotesize
\caption{Độ phân tán mô tả của $T_{work}$ tại một số kích thước.}
\label{tab:dispersion}
\setlength{\tabcolsep}{4pt}
\begin{tabular}{@{}lrrrr@{}}
\toprule
$N$ & Encoding & Trung vị (s) & IQR (s) & IQR/trung vị (\%) \\
\midrule
4 & Binomial & 0.000295 & 0.000004 & 1.4 \\
4 & Binary & 0.000424 & 0.000027 & 6.4 \\
4 & Sequential & 0.000388 & 0.000050 & 13.0 \\
4 & Commander & 0.000500 & 0.000127 & 25.5 \\
4 & Product & 0.000322 & 0.000005 & 1.6 \\
64 & Binomial & 0.5004 & 0.0153 & 3.1 \\
64 & Binary & 0.2273 & 0.009226 & 4.1 \\
64 & Sequential & 3.902 & 0.1286 & 3.3 \\
64 & Commander & 0.3088 & 0.0153 & 4.9 \\
64 & Product & 18.958 & 0.5619 & 3.0 \\
81 & Binomial & 1.091 & 0.0968 & 8.9 \\
81 & Binary & 0.3448 & 0.1239 & 35.9 \\
81 & Sequential & 2.945 & 0.1090 & 3.7 \\
81 & Commander & 0.3738 & 0.0145 & 3.9 \\
81 & Product & 3.648 & 0.3666 & 10.0 \\
216 & Binomial & 22.912 & 0.6474 & 2.8 \\
216 & Binary & 2.296 & 0.0679 & 3.0 \\
216 & Commander & 3.940 & 0.0383 & 1.0 \\
512 & Binary & 81.088 & 7.642 & 9.4 \\
\bottomrule
\end{tabular}
\par\smallskip\begin{minipage}{\linewidth}\footnotesize Q1/Q3 dùng nội suy inclusive trên ba lượt hoàn tất. IQR không phải khoảng tin cậy; không áp dụng bảng này cho nhóm Product N=125 chỉ có một lượt SAT.\end{minipage}
\end{table}

Với ba giá trị đã sắp $t_1\le t_2\le t_3$, trung vị là $t_2$. Nội suy inclusive cho $Q_1=(t_1+t_2)/2$, $Q_3=(t_2+t_3)/2$, nên IQR $=(t_3-t_1)/2$. Script tính trên work của từng lượt, không cộng trung vị riêng của các pha. File \texttt{data/summary\_report.csv} lưu toàn bộ trung vị, Q1/Q3 và trạng thái nhóm. Các cột có tên \texttt{median\_completed} chỉ mô tả các lượt thành công, không được dùng thay cho trung vị của nhóm đủ ba lượt trong bảng chính.

\section{Số literal và ý nghĩa của số đếm}
Binomial, Binary, Sequential và Product đang dùng chỉ có clause AMO hai literal, nên $\ell_E(k)=2c_E(k)$. Commander có điều kiện dừng $\ell_C(k)=2\binom k2$ khi $k\le3$. Với $k>3$:
\[
\ell_C(k)=2I(k)+3k+m+\ell_C(m),\qquad m=\lceil k/3\rceil.
\]
Ba hạng ở tầng hiện tại gồm $2I(k)$ literal nội nhóm, $2k$ literal từ đầu vào đến commander và $k+m$ literal của liên kết chiều ngược. Product cũng có thể được viết dưới dạng truy hồi $\ell_P(k)=4k+\ell_P(p)+\ell_P(q)$ khi $k>4$.

ALO đóng góp $N^2(1+d)$ literal, với $d$ là cờ ALO cột. Toàn bàn cờ có
\[
L_E(N)=N^2(1+d)+(2N+2)\ell_E(N)+4\sum_{k=2}^{N-1}\ell_E(k).
\]
Công thức đúng với các ngưỡng và clause cụ thể đã mô tả. Thay nhóm, ngưỡng dừng hoặc liên kết phải cập nhật công thức và kiểm tra lại với generator. Tổng biến được tính là biến chính cộng biến phụ đã cấp; ID lớn nhất chỉ bằng số biến khi cách cấp ID liên tục và các biến được tính đều xuất hiện.

\section{Đối chiếu yêu cầu và khả năng tái lập}
\begin{table}[pos=H]\centering\small
\caption{Các nội dung đáp ứng yêu cầu Bài tập lớn số 1.}
\begin{tabular}{@{}p{.40\linewidth}p{.53\linewidth}@{}}\toprule
Nội dung & Vị trí trong báo cáo\\\midrule
Giải N-Queens bằng SAT/PySAT & Mô hình CNF, năm encoder và quy trình Glucose3 tại Phần~\ref{sec:methodology}.\\
So sánh CPLEX CP, CPLEX MIP, Gurobi, CP-SAT & Mô hình CP/MIP tại Phần~\ref{sec:methodology}; thời gian, coverage và RAM tại Phần~\ref{sec:experiments}.\\
Thực nghiệm kích thước lớn & 18 kích thước đến 1331; mọi cấu hình đều được thử, không có SKIPPED.\\
Trình bày ALO/AMO & Định nghĩa tại Background; CNF, ví dụ, tính đúng, số đếm tại Methodology.\\
Hình thức bài báo khoa học & Lớp \texttt{cas-sc}; abstract, năm phần chính, Data Availability và tài liệu tham khảo.\\
Liên kết mã nguồn và kết quả & Mã tái tạo tại Phụ lục~\ref{app:reproduce}; URL repository và commit/tag còn cần tác giả điền.\\
\bottomrule\end{tabular}
\end{table}
Các thông tin chưa có trong đầu vào được để rõ để hoàn thiện trước khi nộp: mã sinh viên, lớp/mã học phần, trường/khoa, xác nhận thông tin giảng viên, URL repository và commit/tag mã nguồn. Để tái lập sát môi trường hơn, cần bổ sung tên thương mại CPU, phiên bản executable CP Optimizer và chế độ nguồn điện khi chạy. Không còn phần thực nghiệm đệ quy để trống: Commander và Product trong toàn bộ bảng hiện tại đã là đệ quy.

Tệp manifest lưu fingerprint mã nguồn và cấu hình; script tổng hợp lưu SHA-256 của từng tệp đầu vào trong \texttt{data\_audit.json}. Các dấu vân tay hỗ trợ phát hiện trộn dữ liệu hoặc thay đầu vào, nhưng không thay thế việc công bố mã nguồn. Script cố ý kiểm tra đúng đợt và đúng cấu hình hiện tại; nếu chạy một thí nghiệm khác, cần tạo báo cáo riêng và cập nhật mô tả, không chỉ thay CSV rồi giữ nguyên diễn giải.


\FloatBarrier
\section{Quy trình tái tạo bảng và biểu đồ}
\label{app:reproduce}
Phần hậu xử lý độc lập với các bộ giải. Nó chỉ đọc kết quả đã có và có thể chạy lại nhiều lần để sửa cách trình bày mà không chạy lại thực nghiệm. Tại gốc dự án, cài phụ thuộc trong \texttt{requirements-analysis.txt}, rồi chạy:
\begin{quote}\small\ttfamily
python -m analysis --input results/experiment\_02 --output report
\end{quote}
Đầu vào phải thuộc đợt và biến thể được phân tích; kiểm tra experiment ID/cấu hình giúp ngăn trộn kết quả một tầng với bản đệ quy. Để tái tạo từ bản sao trong gói, dùng \texttt{--input report/data}. Với gói LaTeX độc lập, lệnh \texttt{python build\_report.py} gọi cùng bộ xử lý trên \texttt{data/}.

\begin{table}[pos=H]\centering\small
\caption{Các bước từ dữ liệu gốc đến bảng và hình.}
\begin{tabular}{@{}p{.25\linewidth}p{.68\linewidth}@{}}\toprule
Thành phần & Vai trò và đầu ra\\\midrule
\path{analysis/__main__.py} & Đọc tham số, gọi kiểm tra, tổng hợp, xuất bảng và vẽ; không gọi solver.\\
\path{data_processing.py} & Đối chiếu bốn đầu vào, nhóm theo $N$--phương pháp, tính trung vị/IQR, số đếm và mức độ hoàn tất.\\
\path{tables.py} & Tạo 24 tệp bảng thực nghiệm tại \texttt{report/tables/}; giữ trạng thái thay cho số chưa biết.\\
\path{plots.py} & Tạo tám hình PDF/PNG cùng CSV điểm hình tại \path{report/data/figure_points/}.\\
\path{tests/test_report_analysis.py} & Sáu kiểm thử cho tính đúng của hậu xử lý và khả năng truy vết.\\
\path{report/main.tex} & Ghép văn bản, bảng và hình; biên dịch bằng XeLaTeX.\\
\bottomrule\end{tabular}
\end{table}

\begin{table}[pos=H]\centering\small
\caption{Câu hỏi và quy tắc dữ liệu của tám hình (tên tệp không kèm phần mở rộng).}
\begin{tabular}{@{}p{.25\linewidth}p{.68\linewidth}@{}}\toprule
Tên hình & Mục đích và cách đọc\\\midrule
\texttt{formula\_sizes} & CNF tăng ra sao theo $N$? Vẽ tổng biến và clause đầy đủ theo công thức, kể cả nhóm timeout; không phải tiến độ đã sinh.\\
\texttt{sat\_runtime} & Encoding SAT nào nhanh hơn? Vẽ median work và total; chỉ nhóm hoàn tất 3/3.\\
\texttt{native\_runtime} & SAT so với CP/ILP thế nào? Cùng quy tắc thời gian; so sánh đồng thời mô hình và backend.\\
\texttt{memory} & RSS quan sát tăng thế nào? Median đỉnh RSS cả ba lượt; dấu $\times$ chỉ nhóm chưa hoàn tất 3/3.\\
\texttt{status\_matrix} & Phương pháp giải được những $N$ nào? Mỗi ô giữ đủ trạng thái của ba lượt, gồm UNKNOWN và thành công một phần.\\
\texttt{cactus} & Bao nhiêu kích thước hoàn tất dưới một ngưỡng median total? Trục ngang là thời gian, trục dọc đếm kích thước hoàn tất 3/3.\\
\texttt{sat\_iqr} & Ba lượt dao động ra sao? Median và khoảng Q1--Q3 của work ở nhóm đủ 3/3; không phải khoảng tin cậy.\\
\texttt{phase\_breakdown} & Total lớn do pha nào? Chồng pha của một lượt thật có total trung vị tại $N=4,216$.\\
\bottomrule\end{tabular}
\end{table}

Để kiểm tra một điểm, chạy \texttt{python -m analysis} với \texttt{--input results/experiment\_02 --inspect 64 Binary}. Chương trình in ba lượt và thống kê của nhóm mà không ghi tệp. Đối chiếu trường \texttt{median\_work\_s} với \texttt{summary\_report.csv}, rồi đối chiếu hàng $N=64$, phương pháp Binary, chỉ số median work trong \texttt{figure\_points/sat\_runtime.csv}. Để vẽ riêng hình thời gian SAT, thêm \texttt{--figure sat\_runtime} vào lệnh tạo báo cáo. Bỏ tùy chọn này để tạo toàn bộ bảng và hình.

Các trục log giúp biểu diễn đồng thời nhiều bậc độ lớn: khoảng cách bằng nhau ứng với tỷ số bằng nhau, không phải chênh lệch tuyệt đối bằng nhau. Giá trị thiếu không được thay bằng 0 để vẽ. Đường nối các điểm chỉ hỗ trợ quan sát các kích thước đã thử, không khẳng định thời gian cho các $N$ chưa thực nghiệm. Hình cactus không giữ trục $N$; mỗi bước tăng là một kích thước hoàn tất có median total đạt ngưỡng tương ứng, không phải thời gian tích lũy chạy tất cả bài toán.

Với hình phân rã pha, trường \texttt{coordination} là phần dư giữa total và tổng các pha đã ghi. Đây là giá trị suy ra, không phải số đo trực tiếp một pha mới. Native không có pha load riêng; phần load bằng 0 khi chồng cột chỉ là cách biểu diễn ``không áp dụng''. CSV lưu cờ \texttt{measured}, run ID và số lần lặp để phân biệt. Bảng thành phần có thể trình bày trung vị từng cột, nhưng hình chồng bắt buộc dùng các pha của cùng một lượt để tổng khớp total.

Để kiểm tra mã hậu xử lý, chạy \texttt{python -m unittest discover} với \texttt{-s tests -p test\_report\_analysis.py -v} từ gốc repository sau khi đã tạo đủ hình. Các kiểm thử không gọi lại solver. Môi trường tạo báo cáo được ghi trong \texttt{report\_build\_environment.json}, tách khỏi môi trường chạy solver trong manifest. Khác phiên bản Matplotlib có thể làm thay đổi nhỏ bố cục hình, trong khi các giá trị CSV cần được giữ nguyên.

Repository cần chứa đồng thời mã solver đúng đợt, bốn đầu vào, module \texttt{analysis/}, hướng dẫn, kiểm thử và nguồn báo cáo. Tệp khóa chạy, môi trường ảo và giấy phép thương mại không phải thành phần dữ liệu công bố. Hướng dẫn PowerShell từng bước, cách sửa hình và danh sách tệp cần commit nằm trong \texttt{HUONG\_DAN\_TAI\_TAO.md}. Không có thao tác push GitHub tự động trong bộ mã phân tích.

\FloatBarrier
\begin{thebibliography}{10}
\bibitem{assignment} Học phần Các vấn đề hiện đại KHMT. \emph{Yêu cầu Bài tập lớn số 1}, học kỳ I năm học 2026--2027. Tài liệu đề bài được cung cấp.
\bibitem{courseTemplate} Phạm Quang Minh, Đào Xuân Nghĩa, Tô Văn Khanh. \emph{Exact SAT Solving for the Two-Dimensional Bandwidth Minimization Problem}. Bản thảo dùng làm mẫu báo cáo học phần, 2026.
\bibitem{bach2024} Trần Ngọc Bách. \emph{An optimized SAT approach solving N-Queens problem}. Technology Project Report, University of Engineering and Technology, VNU, 2024.
\bibitem{amoSlides} Tô Văn Khanh. \emph{At Most One -- AMO, At Least One -- ALO Encoding}. Slide Module 3, trang 3--16; tài liệu học phần được cung cấp.
\bibitem{sinz2005} Carsten Sinz. Towards an Optimal CNF Encoding of Boolean Cardinality Constraints. \emph{Principles and Practice of Constraint Programming -- CP 2005}, 827--831. \url{https://www.carstensinz.de/papers/CP-2005.pdf}.
\bibitem{klieber2007} Will Klieber, Gihwon Kwon. \emph{Efficient CNF Encoding for Selecting 1 from N Objects}. 2007. \url{https://www.cs.cmu.edu/~wklieber/papers/2007_efficient-cnf-encoding-for-selecting-1.pdf}.
\end{thebibliography}
